\section{Separate cost models for each graph}
\label{app:type-specialists}
This experiment retains the seven denominator definitions and profile
formulas in Appendix~\ref{app:workflow}. Development uses profile $h=61$,
training uses $h=62,63$, validation uses $h=64$, and the reserved test uses
$h=65,66$. The dimension remains symbolic. Related target batches share
the same split, so their rational kinematic values cannot pass from the
training set into the reserved set. The graph structures and target
patterns occur in both sets.

Let $P$ be the number of physical propagators and let
$\boldsymbol b=(1,\ldots,1,0,\ldots,0)$ have $P$ entries equal to one.
Let $\boldsymbol e_j$ be the unit vector for index $j$. The first batch
contains $\boldsymbol b+2\boldsymbol e_j$ for
$j=1,\ldots,\min(P,3)$. Thus one denominator has power three in each target.
For the second batch, define the lower-sector corner
\begin{equation}
 \boldsymbol b^{\rm low}=
 \begin{cases}
 \boldsymbol b-\boldsymbol e_1,&P<5,\\
 \boldsymbol b-\boldsymbol e_1-\boldsymbol e_3,&P\geq5.
 \end{cases}
\end{equation}
Include this corner and up to three single-dot targets,
$\boldsymbol b^{\rm low}+\boldsymbol e_j$ with $b^{\rm low}_j=1$,
taking the first available indices in increasing order.
The third batch uses the original corner $\boldsymbol b$ and contains
$\boldsymbol b-2\boldsymbol e_j$ for the first three auxiliary entries,
or all auxiliary entries when fewer than three exist. When a family has no auxiliary entry, this construction
repeats the first batch and we remove the duplicate. These formulas define
the inputs; they do not infer their coefficients from the masters.

\subsection{The allowed changes to Kira}
For target batch $\mathcal T$ and a target $\boldsymbol a$, let $S(\boldsymbol a)=\{i:a_i>0\}$ be its
positive-index set. One strategy supplies only maximal target sectors:
\begin{equation}
 \mathcal S_{\rm max}
 =\{S(\boldsymbol a):\boldsymbol a\in\mathcal T,
 \ S(\boldsymbol a)\not\subsetneq S(\boldsymbol a')
 \text{ for every }\boldsymbol a'\in\mathcal T\}.
 \label{eq:type-sectors}
\end{equation}
Kira then includes their subsectors. To see why this is a useful candidate,
consider the shifted integrals produced by the IBP derivative. A raised
index $a_i+1$ carries a factor $a_i$, and the other shifts lower an index.
If $a_i=0$, the raising term vanishes; if $a_i<0$, raising it once leaves
it nonpositive. Thus an IBP equation seeded in one sector stays in that
sector or its subsectors. This observation does not prove that the
restricted pool contains every relation useful for a chosen master basis.
We therefore retain the basis check and full-sector fallback.

Another strategy uses Kira's numerator truncation~\cite{Kira3}. In a sector
with $t$ positive indices, Kira restricts the inherited numerator bound by
\begin{equation}
 s\leq\max(1,t-\tau+1,s_{\rm explicit}),
 \label{eq:type-truncation}
\end{equation}
where $\tau$ is the chosen integer cutoff and $s_{\rm explicit}$ is a
bound explicitly supplied for that sector. We test
$\tau=P-s_c$ and $\tau=P-s_c+1$, where $s_c$ is the largest target numerator
rank. The top-level sector retains its explicit target bound. The first
cutoff combined with maximal target sectors forms the deterministic
control. The development probes also test $\tau=P-s_c-1$.

The remaining choices use the original bounds $B_c$, the smaller bounds
$B_c-(1,0,1)$ with negative entries replaced by zero, the larger bounds
$B_c+(1,1,1)$, Kira integral ordering five~\cite{Kira12}, or the preceding
trained equation filter. Together these choices give nine strategies.
Before each attempt, every strategy enforces a nonempty seed region in
its declared top-level sectors. If $\mathcal S_0$ is that sector list and
$r_{\rm req}$ is the requested positive-index bound, the actual bound is
\begin{equation}
 r_{\rm run}=\max\left(r_{\rm req},\max_{S\in\mathcal S_0}|S|\right).
 \label{eq:type-valid-bound}
\end{equation}
Each positive index is at least one, so a seed in sector $S$ has
$\sum_{i\in S}a_i\geq|S|$. This proves the lower bound used in
Eq.~\eqref{eq:type-valid-bound}. The correction prevents an invalid
configuration from adding a failed startup to any method's cost.
Identical input configurations share one training or validation measurement.
If a proposed reduction fails, it retries with native choices at $B_c$
and then at $B_c+(1,1,1)$. The recorded cost includes every attempt.

\subsection{Fitting and selecting the models}
Each graph has its own predictors of
$y_{cj}=\log(t_{cj}/t_{c0})$, where $t_{cj}$ is the median complete cost of
strategy $j$ on batch $c$ and zero denotes native choices. The nine input
features are the target count; maximum dot count and numerator rank;
minimum and maximum positive-index counts; mean dot count and numerator
rank; number of maximal target sectors; and number of physical
propagators. These features use the requested indices and family
definition. They contain no reduced coefficients.

We fit ten candidates per graph: one constant predictor, trees of maximum
depth one, two and three, a bagged ensemble of sixteen depth-three trees,
ridge regressions with penalties $0.1,1,10$, and nearest-neighbour
predictors with one and three neighbours. Trees use squared-error splits
with at least one row in each child. Each ensemble tree receives a
bootstrap sample of the training batches and uses all features at each
split. This follows bootstrap aggregation~\cite{Breiman1996Bagging}.
For fitted trees $f_1,\ldots,f_{16}$, the ensemble prediction is
\begin{equation}
 \widehat y_j(\boldsymbol f)=\frac1{16}
 \sum_{b=1}^{16}f_{bj}(\boldsymbol f).
\end{equation}
This average also follows by minimizing
$\sum_b(f_{bj}(\boldsymbol f)-v_j)^2$ with respect to $v_j$, using the
leaf-mean derivative in Appendix~\ref{app:workflow}. That appendix derives
the ridge fit and the local-mean estimates used here. Neighbour distances
use the training standard deviations; a constant feature uses divisor
one. We use two timing repetitions for training and three for validation.
Within each batch comparison, every strategy uses the same computer.
Different training batches can use different computers, and the
training response normalizes each cost to its own batch baseline.

Validation considers decision margins $0,0.05,0.15$ in
Eq.~\eqref{eq:workflow-choice}. It chooses the candidate and margin with
the lowest complete cost separately for each graph. The controls are
native choices, the deterministic sector-and-truncation rule, and the
single strategy with the lowest geometric mean training cost for that
graph. A fixed strategy can include the older equation filter. We enable
the new model only when its validation cost ratio to the fastest control
is at most $0.95$. Otherwise, the route uses that control. This threshold
is a selection rule, not a significance test. Unknown graph structures
and features outside the fitted range use native choices.

The reserved comparison uses five repetitions of native choices, the
deterministic rule, the fixed strategy per graph, and the selected route.
The order rotates across batches and repetitions, giving every method
the same number of appearances in each run position over the full test.
Each run uses the same targets, prescribed basis, executable, and one
processor. All methods retain the preceding exact implementation
improvements, so the native control uses that same implementation with
native selection rules. Full time includes input preparation, prediction, native IBP
generation, exact reduction, coefficient checks, cleanup, and retries.
Offline fitting, reference discovery, and saving measurement records stay
outside the timer. Separate application checks receive the basis without
target coefficients and compare their answers only after completion.
The calling program and fitted parameters are loaded before the timing
trials, and each trial starts a fresh Kira process. The additional
application checks also record time with a newly started calling program;
they do not supply a matched speed comparison for that condition.

We combine five-run median ratios within each graph and then give the
seven graphs equal weight, as derived in Appendix~\ref{app:workflow}.
The aggregate intervals resample graph groups. Per-graph plot bars show
the range across the two reserved kinematic profiles. Two profiles do not
support a general significance claim for that integral type, so those
bars are reported as ranges.
